% Authority: exact-scope leaf 15; full PDF page 94; printed p.77.
% U:p077.r28.par.009 | continuation; R28 par.009 visually replayed
\par\noindent
It will suffice to seek a number satisfying the congruence $\chi^2\equiv p\md{b}$; having found
such a number $\chi$, one will always have:
% U:p077.r28.eq.012 | continuation; R28 eq.012 visually replayed
\[
\leg{t}{b}=\leg{\chi(\chi+\psi)}{b}.
\]
% U:p077.r28.par.010 | paragraph-start; R28 par.010 visually replayed
\par\indent
Let us now come to the case in which $\varphi$ is divisible by $b$. The relations $(\delta')$
which then hold are:
% U:p077.r28.eq.013 | continuation; R28 eq.013 visually replayed
\[
t\pm\psi\equiv0\md{b},\qquad
\leg{t\mp\psi}{b}=1.
\]
In the last of these relations one may add to $t\mp\psi$ any multiple of $b$. If one adds to it
the number $t\pm\psi$, which, by virtue of the first relation, is divisible by $b$, there will
result $\leg{2t}{b}=1$, a result which entails this other one: $\leg{t}{b}=\leg{2}{b}$.
% U:p077.r28.par.011 | paragraph-start; R28 par.011 visually replayed
\par\indent
Combining this result with a known theorem, one will see that, in the case of $\varphi$ divisible
by $b$, one has $\leg{t}{b}=1$ or $\leg{t}{b}=-1$, according as $t$ is of the form $8n+7$ or of
this one: $8n+3$.
% U:p077.r28.par.012 | paragraph-start; R28 par.012 visually replayed
\par\indent
Comparing what has just been proved with the theorem established at the end of the preceding
paragraph, one will see that it is possible to decide, independently of the knowledge of the
number $t$, whether $-b$ is or is not a biquadratic residue with respect to $p$. Since the result
thus reached no longer contains any trace of the number $t$, one is led to believe that it does
not presuppose the possibility of the equation $t^2-bu^2=p$, from which the number $t$ takes its
origin, and that it is generally true for all values of $b$ and $p$ such that $\leg{b}{p}=1$. This
is indeed what takes place, as one can prove, by examining, instead of the equation $t^2-bu^2=p$,
the more general equation $t^2\pm Mu^2=p$, where $M$ denotes the product of any number of distinct
prime numbers, and then combining the results of this examination with the following theorem, of
whose truth one cannot doubt, but whose rigorous demonstration nevertheless presents difficulties:
% U:p077.r28.par.013 | paragraph-start; R28 par.013 visually replayed
\par\indent
“$c$ denoting any prime number whatever and $p$ a prime number $4n+1$, such that $\leg{c}{p}=1$,
one can always determine a number $\delta$, composed of simple factors all smaller than $c$, and
such that the equation $t^2\pm c\delta u^2=p$, is soluble.”
% U:p077.r28.par.014a | paragraph-start; R28 par.014a visually replayed
\par\indent
However that may be, one may remark that the theorem which we are about to state is rigorously
proved, by what precedes, for all values
